\section{Modèle du système et formulation du problème}

Cette section traduit le scénario présenté dans l'introduction en un modèle
opérationnel et coopératif précis. Elle introduit les opérateurs, leurs
demandes, leurs capacités et leurs coûts, caractérise les décisions
admissibles, puis définit le jeu des économies et son critère de stabilité.

L'accord finalement recherché réunit tous les opérateurs présents sur le site.
Le même problème opérationnel doit néanmoins être défini pour chaque groupe
susceptible de coopérer seul. Il devient ainsi possible de comparer la
coopération globale aux options de repli dont disposent ses membres. Les
notations nécessaires à cette comparaison sont introduites ci-dessous.


\subsection{Cadre et hypothèses}

Nous considérons un ensemble d'opérateurs de réseaux mobiles
\(N=\{1,\dots,n\}\), avec \(n\geqslant2\), présents sur un même site radio et
desservant une même zone géographique. Chaque opérateur \(i\in N\) dispose
d'un équipement radio agrégé qui peut être maintenu actif ou placé en veille.

L'horizon fondamental est une fenêtre \(H\), finie, non vide et composée
d'heures consécutives. La demande de l'opérateur \(i\) à l'heure \(h\in H\)
est notée \(d_i^h\). Afin
d'introduire les briques opérationnelles sans alourdir les notations, nous
fixons d'abord une heure arbitraire \(h\in H\) et écrivons
\(d_i:=d_i^h\). Les objets sans indice \(h\) définis ci-dessous sont donc
horaires ; ils ne constituent pas un second modèle. La décision sur la fenêtre
est ensuite obtenue en les articulant sous une contrainte commune de gardiens.
Dans les hypothèses suivantes, un groupe coopérant est simplement appelé
\emph{coalition} ; sa définition formelle et les décisions qui lui sont
associées sont introduites dans les sous-sections suivantes.

\begin{enumerate}
    \item \textbf{Site unique et recouvrement de couverture.}
    Les équipements considérés couvrent la même zone. Tout sous-ensemble
    non vide d'équipements actifs est donc supposé assurer la couverture
    du site. L'extinction d'un équipement ne modifie ni la couverture
    utile ni les interférences environnantes. La faisabilité d'une
    configuration de gardiens dépend donc uniquement de la capacité
    effective agrégée de leurs équipements à acheminer le trafic.

    \item \textbf{Fongibilité du trafic, qualité de service et capacités
    effectives.}
    Le trafic des membres d'une coalition peut être acheminé
    indifféremment par chacun de ses opérateurs actifs. Les mécanismes
    techniques nécessaires à ce transfert sont supposés disponibles, et le
    transfert ne modifie pas la qualité de service perçue. Le débit, la latence
    et la disponibilité ne sont pas modélisés séparément : leurs exigences
    sont intégrées aux capacités effectives, et toute répartition admissible
    est supposée les respecter. Ces capacités sont indépendantes de l'identité
    des clients servis, additives, et toute la demande doit être satisfaite.

    \item \textbf{Information complète.}
    Les demandes \((d_i^h)_{i\in N,h\in H}\), les capacités et les paramètres
    de coût sont connus sur toute la fenêtre. Le modèle ne comporte ni
    incertitude ni variable d'état décrivant l'historique d'utilisation des
    équipements.

    \item \textbf{États opérationnels binaires.}
    Chaque équipement est soit actif et susceptible d'acheminer du
    trafic, soit en veille et n'en achemine aucun. Les délais et coûts de
    transition entre ces états sont négligés, ce qui correspond à une
    heure de décision suffisamment longue devant les temps de
    commutation. Seuls les coûts opérationnels évitables sont modélisés ;
    les coûts structurels supportés indépendamment de l'état de
    l'équipement sont exclus. Dans cette étude, l'état de veille est supposé
    consommer \(0\ \mathrm{W}\).

    \item \textbf{Coût opérationnel affine en taux de charge.}
    Lorsqu'un équipement est actif, son coût opérationnel évitable est
    la somme d'une composante fixe et d'une composante proportionnelle à
    son taux de charge. Cette structure reprend les modèles de puissance
    séparant une consommation active à vide d'une contribution dépendant de
    la charge \cite{arnold2010,auer2011,lorincz2012}. La composante fixe
    représente la consommation active à vide ; la composante variable,
    l'accroissement de consommation lié aux ressources mobilisées.

    Le volume descendant horaire est utilisé comme proxy de la charge servie,
    et les coefficients sont estimés à partir des couples mesurés de trafic et
    de puissance. L'hypothèse est que, sur le domaine observé, leur relation
    agrégée est affine au premier ordre ; sa validité empirique est évaluée
    dans la Section~\ref{sec:simulations_numeriques}.

    \item \textbf{Neutralité des revenus et utilité transférable.}
    Toute solution opérationnelle admissible sert la même demande totale. Les
    revenus tirés des abonnés sont donc supposés indépendants de
    l'équipement qui achemine physiquement leur trafic et n'interviennent
    pas dans la comparaison des solutions opérationnelles. Les économies
    opérationnelles peuvent être redistribuées sans friction entre les
    membres d'une coalition au moyen de transferts monétaires.
\end{enumerate}

La décision opérationnelle comporte ainsi deux niveaux : le choix des
équipements actifs et la répartition du trafic entre eux. Les
paramètres associés sont introduits ci-dessous.

\subsection{Opérateurs et fonctionnement autonome}

Pour l'heure arbitraire fixée ci-dessus, toutes les grandeurs de trafic
désignent des volumes cumulés sur une heure et sont exprimées en gigaoctets
(\(\mathrm{Go}\)). Toutes les grandeurs monétaires désignent les coûts
supportés pendant cette heure et sont exprimées en euros (\(\text{€}\)).

\begin{definition}[Paramètres individuels et coût opérationnel autonome]
Pour tout opérateur \(i\in N\), les paramètres primitifs du modèle sont :
\begin{itemize}
    \item \textbf{Demande \(d_i > 0\).}
    Quantité de trafic engendrée par les clients de l'opérateur \(i\)
    pendant l'heure considérée, exprimée en \(\mathrm{Go}\).

    \item \textbf{Capacité effective \(q_i>0\).}
    Quantité maximale de trafic que l'équipement actif de l'opérateur
    \(i\) peut acheminer pendant cette heure, également exprimée en
    \(\mathrm{Go}\).

    \item \textbf{Coût fixe de fonctionnement en état actif \(F_i>0\).}
    Coût énergétique évitable supporté pendant l'heure considérée dès
    que l'équipement de l'opérateur \(i\) est actif, indépendamment du
    trafic acheminé. Il
    correspond à la valorisation monétaire de la consommation électrique
    incompressible en fonctionnement. Le paramètre \(F_i\) est exprimé en
    \(\text{€}\).

    \item \textbf{Surcoût variable à pleine charge \(\beta_i>0\).}
    Surcoût énergétique supporté pendant l'heure considérée lorsque le taux
    de charge de l'équipement vaut un. Il correspond à la valorisation
    monétaire de la consommation électrique additionnelle liée au trafic.
    Le paramètre \(\beta_i\) est exprimé en \(\text{€}\) ; il ne constitue pas
    un coût variable unitaire par \(\mathrm{Go}\).
\end{itemize}

\paragraph{Périmètre du coût.}
Dans tout l'article, le coût opérationnel désigne la valorisation monétaire
de l'énergie électrique consommée pendant une heure. Sous un
prix de l'électricité commun et strictement positif, le passage de l'énergie
au coût est une multiplication par une constante positive. La minimisation
du coût, ses comparaisons entre coalitions et les économies qui en découlent
sont donc exactement équivalentes à leurs contreparties énergétiques.

Seuls les coûts énergétiques qui dépendent de la décision opérationnelle
sont retenus. Les charges non énergétiques et les consommations supportées
indépendamment de l'état de l'équipement sont omises. Le coût de la veille
est nul, tandis que celui de l'état actif au taux de charge
\(\rho\in[0,1]\) vaut \(F_i+\beta_i\rho\).
Sous l'hypothèse de veille nulle, \(F_i\) représente directement la composante
fixe évitable de l'état actif. Les coûts structurels indépendants de l'état
restent exclus.

Le fonctionnement autonome est supposé réalisable :
\[
    d_i\leqslant q_i,\qquad \forall i\in N.
\]

Les grandeurs dérivées associées à l'opérateur \(i\) sont les suivantes :
\begin{itemize}
    \item \textbf{Taux de charge autonome \(\rho_i^0\).}
    \[
        \rho_i^0:=\frac{d_i}{q_i}\in(0,1].
    \]
    Cette grandeur sans dimension mesure la fraction de la capacité
    utilisée lorsque l'opérateur sert seul sa demande.

    \item \textbf{Coût variable unitaire \(\gamma_i\).}
    \[
        \gamma_i:=\frac{\beta_i}{q_i}.
    \]
    Il représente l'accroissement du coût opérationnel par \(\mathrm{Go}\)
    supplémentaire de trafic. Il est exprimé en
    \(\text{€}/\mathrm{Go}\). Ainsi, pour un trafic \(t_i\), le coût variable
    \(\beta_i(t_i/q_i)\) s'écrit également \(\gamma_i t_i\).

    \item \textbf{Coût opérationnel autonome \(C_i^0\).}
    \[
        C_i^0:=F_i+\beta_i\rho_i^0
        =F_i+\gamma_i d_i.
    \]
    Il s'agit du coût évitable supporté par l'opérateur \(i\) lorsqu'il
    maintient son équipement actif et achemine sa propre demande.
\end{itemize}
\end{definition}



\subsection{Coalitions, configurations de gardiens et répartitions admissibles}


\begin{definition}[Coalition]
On note \(2^N:=\{A\mid A\subseteq N\}\) l'ensemble des sous-ensembles de
\(N\), et \(\mathcal S:=2^N\setminus\{\varnothing\}\) l'ensemble des
coalitions. Une \textbf{coalition} est donc un élément \(S\in\mathcal S\) ;
elle représente un groupe d'opérateurs coopérant sur le site. Ainsi,
\(|\mathcal S|=2^n-1\).
\end{definition}

Pour tout $A\subseteq N$, sa \textbf{demande de trafic totale} et sa
\textbf{capacité effective agrégée} sont respectivement
\[
    d(A):=\sum_{i\in A}d_i,
    \qquad
    q(A):=\sum_{i\in A}q_i.
\]

\begin{definition}[Configuration de gardiens]
Dans une coalition $S \in \mathcal S$, un ensemble d'\textbf{opérateurs
gardiens} est un \(G\in\mathcal S\) tel que \(G\subseteq S\), dont les
équipements restent actifs et acheminent la totalité de la demande de trafic
\(d(S)\).

Le couple $(S,G)$ est appelé \textbf{configuration de gardiens}. Il est dit \textbf{faisable} si la capacité effective agrégée des équipements des gardiens est suffisante pour acheminer le trafic de la coalition, soit $d(S) \leqslant q(G)$.
Pour tout $S \in \mathcal S$, on note $\mathcal G(S)$ l'ensemble des
configurations de gardiens faisables :
\[
    \mathcal G(S)
    :=\left\{G\in\mathcal S\mid G\subseteq S,\ d(S)\leqslant q(G)\right\}.
\]

On note $\mathcal{E} = \left\{ (S, G) \;\middle|\; S \in \mathcal S \text{ et } G \in \mathcal G(S) \right\}$ l'ensemble de toutes les configurations de gardiens faisables.

L'hypothèse \(d_i\leqslant q_i\) pour tout \(i\in N\) implique
\(d(S)\leqslant q(S)\). Ainsi, \(S\in\mathcal G(S)\), de sorte que
\(\mathcal G(S)\) est non vide pour tout \(S\in\mathcal S\), et donc que
\(\mathcal E\) est non vide.

\end{definition}


\begin{definition}[Vecteur de répartition du trafic admissible]
Soit $(S,G)\in\mathcal E$. Un \textit{vecteur de répartition du trafic}
est un vecteur global $\mathbf t=(t_i)_{i\in N}\in\mathbb R_+^N$, où
$t_i$ désigne le volume de trafic, exprimé en \(\mathrm{Go}\), acheminé par
l'équipement de l'opérateur $i$. L'ensemble des répartitions admissibles
pour $(S,G)$ est
\[
\mathcal T(S,G):=
\left\{
\mathbf t\in\mathbb R_+^N\ \middle|\
\begin{aligned}
    &0\leqslant t_i\leqslant q_i, &&\forall i\in G,\\
    &\sum_{i\in G}t_i=d(S),\\
    &t_i=0, &&\forall i\in N\setminus G
\end{aligned}
\right\}.
\]
\end{definition}

Pour toute répartition \(\mathbf t\in\mathcal T(S,G)\), le
\textbf{taux de charge} de l'équipement de l'opérateur \(i\) est
\[
    \rho_i(\mathbf t):=\frac{t_i}{q_i}\in[0,1].
\]
En particulier, \(\rho_i(\mathbf t)=0\) pour \(i\notin G\), tandis que
\(\rho_i(\mathbf t)=1\) caractérise un équipement saturé.

    
\subsection{Coût opérationnel et solution optimale}


\begin{definition}[Coût d'une solution opérationnelle admissible]
Soit $(S, G)\in\mathcal E$ et $\mathbf{t}\in\mathcal{T}(S, G)$. Le triplet $(S,G,\mathbf t)$ est appelé \textbf{solution opérationnelle admissible}. Son \textbf{coût opérationnel total} est défini par
\[
C(S,G , \mathbf{t})
=\sum_{i\in G}\left(F_i+\beta_i\rho_i(\mathbf t)\right)
=\underbrace{\sum_{i\in G}F_i}_{\text{coûts fixes}}
+\underbrace{\sum_{i\in G}\gamma_i t_i}_{\text{coûts variables}}
\]
\end{definition}

\begin{definition}[Coût minimal d'une coalition]
Soit $S\in \mathcal S$. On définit le coût minimal de $S$ par
\[
C^{*}(S)=\min_{G \in\mathcal{G}(S)}\min_{\mathbf{t}\in\mathcal T(S, G)}C(S, G, \mathbf{t})
\]
c'est-à-dire le plus faible coût opérationnel que la coalition $S$ puisse atteindre en choisissant une configuration de gardiens faisable et une répartition du trafic admissible.
\end{definition}

Pour tout \(G\in\mathcal G(S)\), l'ensemble \(\mathcal T(S,G)\) est un
polytope non vide et compact ; comme \(C(S,G,\cdot)\) est continue et
\(\mathcal G(S)\) est fini, les deux minima intervenant dans la définition
de \(C^*(S)\) sont atteints.


\begin{remarque}
\label{rem:eq}
Pour tout opérateur $i\in N$, l'unique ensemble de gardiens admissible pour la coalition singleton $\{i\}$ est $\{i\}$, et l'unique répartition admissible satisfait $t_i=d_i$. Par conséquent,
\[
C^*(\{i\})=F_i+\gamma_i d_i=C_i^0.
\]
\end{remarque}

\subsection{Fenêtre de décision et politiques opérationnelles}
\label{subsec:fenetre_decision}

Pour tout \(h\in H\) et tout \(S\in\mathcal S\), posons
\[
    d^h(S):=\sum_{i\in S}d_i^h.
\]
La demande cumulée sur la fenêtre est notée
\[
    D_i:=\sum_{h\in H}d_i^h,
    \qquad
    D(S):=\sum_{h\in H}d^h(S)=\sum_{i\in S}D_i.
\]
Le fonctionnement autonome est réalisable à chaque heure :
\[
    d_i^h\leqslant q_i,
    \qquad \forall i\in N,\ \forall h\in H.
\]
Les capacités \(q_i\) et les paramètres de coût horaire
\(F_i,\gamma_i\) sont supposés constants sur \(H\). Ainsi, le coût fixe
\(F_i\) est supporté à chaque heure pendant laquelle l'équipement reste
actif ; il ne s'agit pas d'un coût ponctuel de commutation. À chaque heure,
les définitions précédentes s'appliquent en remplaçant \(d_i\) par
\(d_i^h\). Les ensembles de gardiens, les répartitions admissibles, le coût
et le coût optimal horaires sont respectivement notés
\(\mathcal G_h(S)\), \(\mathcal T_h(S,G)\),
\(C_h(S,G,\mathbf t^h)\) et \(C_h^*(S)\).

La politique principale impose qu'un même ensemble de gardiens reste actif
pendant toute la fenêtre. Pour tout \(S\in\mathcal S\), les ensembles
admissibles sont donc
\[
    \mathcal G_H(S)
    :=
    \left\{
        G\in\mathcal S
        \ \middle|\
        G\subseteq S,\ d^h(S)\leqslant q(G),\ \forall h\in H
    \right\}.
\]
Une solution sur \(H\) est formée d'un ensemble
\(G\in\mathcal G_H(S)\) et d'une répartition
\(\mathbf t^h\in\mathcal T_h(S,G)\) pour chaque heure. Les gardiens sont
fixes sur la fenêtre, mais les volumes qu'ils acheminent peuvent varier avec
la demande. Cette persistance constitue une contrainte d'engagement
opérationnel sur \(H\), et non la conséquence d'un coût de commutation inclus
dans la fonction objectif.

Le coût optimal de cette politique, dite \emph{persistante}, est
\begin{equation}
\begin{aligned}
    C_H^*(S)
    :=
    \min_{G\in\mathcal G_H(S)}
    \min_{(\mathbf t^h)_{h\in H}
    \in\prod_{h\in H}\mathcal T_h(S,G)}
    \sum_{h\in H}C_h(S,G,\mathbf t^h).
    \label{eq:cout_persistant}
\end{aligned}
\end{equation}

Une politique alternative autorise au contraire un nouvel ensemble de
gardiens à chaque heure. Son coût optimal est
\begin{equation}
    C_{H,\mathrm{hor}}^*(S)
    :=
    \sum_{h\in H}C_h^*(S)
    \leqslant C_H^*(S).
    \label{eq:cout_horaire_borne}
\end{equation}
Cette politique \emph{horairement reconfigurable} est un problème
opérationnel distinct, et non une autre interprétation de la politique
persistante. En l'absence de coût de commutation, elle fournit la référence
optimiste à laquelle la politique persistante peut être comparée. Lorsque
\(H\) contient une seule heure, les deux politiques coïncident avec le
problème horaire défini plus haut.

\subsection{Jeu des économies et critère de stabilité}
\label{subsec:modele_jeu_economies}

Le fonctionnement autonome fournit une référence naturelle pour mesurer les
gains de la coopération sur \(H\). Pour un singleton, les deux politiques
coïncident et
\[
    C_H^*(\{i\})
    =C_{H,\mathrm{hor}}^*(\{i\})
    =\sum_{h\in H}\left(F_i+\gamma_i d_i^h\right).
\]

\begin{definition}[Jeux des économies sur la fenêtre]
\label{def:jeu_economies}
Le jeu principal, associé à la politique persistante, est défini par
\begin{equation}
    v_H(S)
    :=
    \sum_{i\in S}C_H^*(\{i\})-C_H^*(S),
    \qquad S\in\mathcal S.
    \label{eq:jeu_persistant}
\end{equation}
Le jeu associé à la politique horairement reconfigurable est
\begin{equation}
    v_{H,\mathrm{hor}}(S)
    :=
    \sum_{i\in S}C_H^*(\{i\})-C_{H,\mathrm{hor}}^*(S),
    \qquad S\in\mathcal S,
    \label{eq:jeu_horaire}
\end{equation}
On pose \(v_H(\varnothing)=v_{H,\mathrm{hor}}(\varnothing)=0\). Ces valeurs
représentent les économies de coût énergétique relativement au fonctionnement
autonome et, à un facteur de conversion positif près, l'énergie évitée.
\end{definition}

La référence autonome étant commune aux deux jeux, la borne
\eqref{eq:cout_horaire_borne} implique, pour tout \(S\in\mathcal S\),
\[
    v_{H,\mathrm{hor}}(S)-v_H(S)
    =
    C_H^*(S)-C_{H,\mathrm{hor}}^*(S)
    \geqslant0.
\]
Cet écart mesure exactement l'économie supplémentaire permise par un changement
de gardiens au sein de la fenêtre.

L'analyse de stabilité de la Section~\ref{sec:jeu_cooperatif} porte sur le
jeu principal \((N,v_H)\). Le jeu \((N,v_{H,\mathrm{hor}})\) sert à isoler
l'effet opérationnel de la possibilité de changer de gardiens au sein de la
fenêtre ; il n'est jamais confondu avec \((N,v_H)\).

Une répartition peut être acceptable pour chaque opérateur pris
isolément, tout en étant rejetée par un groupe d'opérateurs. En effet, une
sous-coalition $Q$ a intérêt à quitter $N$ si ses membres reçoivent
ensemble moins que l'économie $v_H(Q)$ qu'ils pourraient réaliser seuls.
Nous cherchons donc une répartition qui distribue toute l'économie de
$N$ et écarte toute déviation collective.

Pour $z=(z_i)_{i\in N}\in\mathbb R^N$ et $A\subseteq N$, notons
\[
    z(A):=\sum_{i\in A}z_i.
\]

\begin{definition}[Allocation des économies]
Une allocation des économies du jeu $(N,v_H)$ est un vecteur
$z=(z_i)_{i\in N}\in\mathbb R^N$ tel que
\[
    z(N)=v_H(N)
    \qquad\text{et}\qquad
    z_i\geqslant0,
    \quad \forall i \in N.
\]
L'ensemble de ces allocations est noté $\mathcal A(N,v_H)$.
\end{definition}

La condition $z_i\geqslant0$ exprime la rationalité individuelle puisque
$v_H(\{i\})=0$ : aucun opérateur ne doit supporter un coût net supérieur à
son coût autonome.

\begin{definition}[Coeur du jeu des économies]
\label{def:coeur_v}
Le coeur du jeu $(N,v_H)$ est l'ensemble
\[
\mathcal C(N,v_H)
:=
\left\{
z\in\mathcal A(N,v_H)
\;\middle|\;
z(Q)\geqslant v_H(Q),
\quad
\forall\,Q\in\mathcal S
\right\}.
\]
\end{definition}

Pour tout groupe susceptible de faire sécession, la contrainte
$z(Q)\geqslant v_H(Q)$ garantit que ses membres reçoivent collectivement au
moins l'économie qu'ils pourraient obtenir en coopérant séparément.

Le problème comporte désormais deux étapes logiquement distinctes. La
Section suivante calcule les coûts optimaux et établit si la réunion de
coalitions crée au moins autant d'économies que leur fonctionnement séparé.
La Section~\ref{sec:jeu_cooperatif} détermine ensuite si ces économies peuvent
être réparties de manière stable entre tous les opérateurs.
